%% Copyright 2026 Petr Socha
%%
%% This work may be distributed and/or modified under the
%% conditions of the LaTeX Project Public License, either version 1.3c
%% of this license or (at your option) any later version.
%% The latest version of this license is in
%%   https://www.latex-project.org/lppl.txt
%% and version 1.3c is part of all distributions of LaTeX
%% version 2008 or later.
%%
%% This work has the LPPL maintenance status `maintained'.
%% The Current Maintainer of this work is Petr Socha.
%% Maintenance and bug reports: https://github.com/petrsocha/pslib-latex/issues
%%
%% The components of this work are listed in MANIFEST.txt.
%% The school logo is a separately licensed third-party asset.

% ============================================================================
% ANALÝZA -- NADPISY, SEZNAMY A CITACE
% ============================================================================
\chapter{Analýza}
\label{chap:analyza}

% Nadpisy v LaTeXu: \section, \subsection a případně \subsubsection.
% Čísla nadpisů se doplňují automaticky. Nezapisujte je ručně.
\section{Práce s odbornými zdroji}

Základní principy analogových a~digitálních obvodů lze studovat například
v~knize \emph{Foundations of Analog and Digital Electronic Circuits}~\cite{agarwal2005}.

% Citaci vložte pomocí \cite{klic}. Klíč musí existovat v literatura.bib.
% Ukázka: odborná kniha (typ @book). Další druhy zdrojů najdete v ostatních kapitolách.
% Pokud čerpáte konkrétní myšlenku, umístěte citaci k příslušné větě.

\section{Požadavky}
\label{sec:pozadavky}

% Odrážkový seznam začíná prostředím itemize.
\subsection{Funkční požadavky}
Příklady požadavků na jednoduchou aplikaci:
\begin{itemize}
  \item zobrazit naměřenou hodnotu,
  \item uložit výsledek měření,
  \item upozornit na překročení limitu.
\end{itemize}

% Číslovaný seznam se zapisuje do prostředí enumerate.
\subsection{Nefunkční požadavky}
Pro ilustraci lze vymezit i další požadavky:
\begin{enumerate}
  \item snadná obsluha,
  \item přehledné rozhraní,
  \item spolehlivé ukládání dat.
\end{enumerate}

% \autoref přidá automaticky i název typu odkazu („oddíl“, „kapitola“).
% Pokud potřebujete český pád („v oddílu“), napište jej ručně s \ref.
Požadavky shrnuje \autoref{sec:pozadavky}. Na tuto analýzu
navazuje \autoref{chap:navrh}.
